\documentclass[11pt,letterpaper]{article}

\usepackage[utf8]{inputenc}
\usepackage[T1]{fontenc}
\usepackage[spanish,es-tabla]{babel}
\usepackage{lmodern}
\usepackage{geometry}
\usepackage{microtype}
\usepackage{amsmath,amssymb,amsfonts}
\usepackage{booktabs}
\usepackage{longtable}
\usepackage{array}
\usepackage{enumitem}
\usepackage{xcolor}
\usepackage{hyperref}
\usepackage{graphicx}
\usepackage{float}
\usepackage{titlesec}
\usepackage{caption}
\usepackage{fancyhdr}
\usepackage{listings}
\usepackage{url}

\geometry{margin=2.3cm}
\hypersetup{
    colorlinks=true,
    linkcolor=blue!50!black,
    citecolor=blue!50!black,
    urlcolor=blue!50!black,
    pdftitle={Atlas topológico de exoplanetas con Mapper},
    pdfauthor={Equipo de proyecto},
    pdfsubject={Topological Data Analysis aplicado a exoplanetas},
    pdfkeywords={Mapper, TDA, exoplanetas, sesgo observacional, machine learning}
}

\setlength{\headheight}{14pt}
\pagestyle{fancy}
\fancyhf{}
\lhead{Atlas topológico de exoplanetas}
\rhead{Mapper + TDA}
\cfoot{\thepage}

\setlist[itemize]{topsep=2pt,itemsep=2pt,parsep=1pt,leftmargin=1.2cm}
\setlist[enumerate]{topsep=2pt,itemsep=2pt,parsep=1pt,leftmargin=1.2cm}
\renewcommand{\arraystretch}{1.18}

\lstdefinestyle{pythonstyle}{
    basicstyle=\ttfamily\small,
    breaklines=true,
    frame=single,
    rulecolor=\color{black!30},
    backgroundcolor=\color{black!3},
    keywordstyle=\color{blue!60!black},
    commentstyle=\color{green!40!black},
    stringstyle=\color{red!50!black},
    showstringspaces=false
}
\lstset{style=pythonstyle}

\newcommand{\Rp}{R_p}
\newcommand{\Mp}{M_p}
\newcommand{\rhop}{\rho_p}
\newcommand{\Teq}{T_{\rm eq}}
\newcommand{\Mapper}{\operatorname{Mapper}}
\newcommand{\phys}{\rm phys}
\newcommand{\orb}{\rm orb}
\newcommand{\joint}{\rm joint}
\newcommand{\bias}{\rm bias}
\newcommand{\E}{\mathbb{E}}
\newcommand{\R}{\mathbb{R}}

\title{\vspace{-1.0cm}\textbf{Atlas topológico de exoplanetas con Mapper}\\
\large Planteamiento, metodología, hipótesis, validación y ruta hacia machine learning}
\author{Proyecto de análisis topológico de poblaciones exoplanetarias}
\date{\today}

\begin{document}
\maketitle

\begin{abstract}
Este documento plantea un proyecto de análisis topológico de datos para estudiar la población de exoplanetas confirmados usando el algoritmo Mapper. El objetivo es construir y comparar tres representaciones: un Mapper físico-composicional, un Mapper orbital-energético y un Mapper conjunto. El proyecto busca responder si las estructuras topológicas inducidas por variables físicas como radio, masa y densidad se parecen, divergen o se acoplan parcialmente con las estructuras inducidas por variables orbitales como período, semieje mayor, insolación y temperatura de equilibrio. Además, se propone un módulo explícito para detectar sesgo observacional por método de descubrimiento, usando métricas de pureza, entropía, enriquecimiento, asociación global, permutación y modelos supervisados. El resultado esperado no es una sola gráfica, sino un atlas topológico validado: una colección coordinada de mapas, métricas y controles que permitan distinguir estructura astrofísica, transición poblacional y sesgo instrumental.
\end{abstract}

\tableofcontents
\newpage

\section{Glosario técnico}

\begin{longtable}{p{0.25\linewidth}p{0.68\linewidth}}
\toprule
\textbf{Término} & \textbf{Definición operativa} \\
\midrule
\endhead
Exoplaneta & Planeta orbitando una estrella distinta del Sol. En este proyecto se consideran exoplanetas confirmados registrados en NASA Exoplanet Archive. \\
Espacio de características & Matriz $X\in\R^{n\times p}$ donde cada fila representa un planeta y cada columna una variable física, orbital, estelar u observacional. \\
Variable física & Magnitud intrínseca o derivada del planeta: radio, masa, densidad. \\
Variable orbital & Magnitud que describe la órbita o ambiente energético: período, semieje mayor, insolación, temperatura de equilibrio, excentricidad. \\
Variable observacional & Magnitud que describe el modo de detección o selección: método de descubrimiento, año, instalación, distancia, banderas de tránsito o velocidad radial. \\
Mapper & Algoritmo de Topological Data Analysis que transforma datos de alta dimensión en un grafo mediante funciones filtro, cubiertas solapadas y clustering local \cite{singh2007mapper}. \\
Atlas topológico & Colección coordinada de Mappers, métricas y controles que describen el mismo sistema desde varias vistas: física, orbital, conjunta, observacional y predictiva. \\
Sesgo observacional & Distorsión de la muestra causada por limitaciones de detección, instrumentos, geometría observacional o selección del catálogo. \\
Topología de la red & Propiedades estructurales del grafo: componentes, ciclos, ramas, puentes, comunidades, diámetro, centralidad y conectividad espectral. \\
Puente topológico & Región del grafo que conecta dos poblaciones y puede representar transición física, mezcla de clases o artefacto de construcción. \\
\bottomrule
\end{longtable}

\section{Resumen ejecutivo}

El proyecto propone estudiar exoplanetas con tres preguntas centrales:

\begin{enumerate}
    \item \textbf{Estructura física:} ¿cómo se organiza la población de exoplanetas en el espacio de radio, masa y densidad?
    \item \textbf{Estructura orbital:} ¿cómo se organiza la población en el espacio de período, semieje mayor, insolación y temperatura de equilibrio?
    \item \textbf{Acoplamiento físico-orbital:} ¿las estructuras de ambos espacios son similares, independientes o parcialmente acopladas?
\end{enumerate}

La estrategia consiste en construir tres grafos Mapper:

\begin{equation}
G_{\phys}=\Mapper(X_{\phys}),\qquad
G_{\orb}=\Mapper(X_{\orb}),\qquad
G_{\joint}=\Mapper(X_{\joint}).
\end{equation}

Luego se comparan sus estructuras mediante métricas de grafos, distancias topológicas, etiquetas físicas y controles observacionales. El producto final debe ser un \textbf{atlas topológico de exoplanetas}: no una visualización aislada, sino un sistema reproducible de mapas, métricas, pruebas y conclusiones.

\section{Contexto y justificación}

La base de datos propuesta es la tabla \texttt{PSCompPars} del NASA Exoplanet Archive. Esta tabla compila parámetros planetarios, estelares y de sistema para exoplanetas confirmados y está diseñada para análisis estadístico de poblaciones, con una fila por planeta \cite{nasa_pscp_about,nasa_ps_columns}. Sin embargo, la propia documentación advierte que los valores de una fila pueden venir de distintas fuentes o ser calculados; por tanto, pueden no ser internamente autoconsistentes \cite{nasa_pscp_calc}. Esta característica obliga a diseñar controles de calidad, filtrado, imputación y análisis de sensibilidad.

La motivación astrofísica es que las poblaciones exoplanetarias no parecen distribuirse de forma uniforme. Por ejemplo, la literatura reporta una deficiencia de planetas pequeños entre aproximadamente $1.5$ y $2.0\,R_\oplus$, conocida como \emph{radius valley} o \emph{Fulton gap} \cite{fulton2017gap}. También se han reportado regiones escasas de planetas tipo Neptuno en períodos cortos, conocidas como \emph{Neptunian desert} o \emph{sub-Jovian desert} \cite{mazeh2016desert}. Estas estructuras pueden aparecer como huecos, ramas, puentes o separaciones en representaciones topológicas.

El algoritmo Mapper es apropiado porque permite estudiar datos de alta dimensión como una red interpretable. Mapper combina tres ideas: filtros, cubiertas solapadas y clustering local. Su salida no es una clasificación rígida, sino una representación de la forma global de los datos \cite{singh2007mapper}.

\section{Planteamiento del problema}

\subsection{Problema general}

Queremos saber si la estructura física de los exoplanetas y la estructura orbital de los exoplanetas generan topologías compatibles, divergentes o parcialmente acopladas.

En forma compacta:

\begin{equation}
\mathcal{T}(X_{\phys}) \overset{?}{\sim} \mathcal{T}(X_{\orb}) \overset{?}{\subseteq} \mathcal{T}(X_{\joint}),
\end{equation}

donde $\mathcal{T}(X)$ representa la estructura topológica extraída de $X$ mediante Mapper.

\subsection{Problemas específicos}

\begin{enumerate}
    \item Identificar regiones físicas: terrestres, super-Tierras, sub-Neptunos, gigantes gaseosos y super-Júpiteres.
    \item Identificar regiones orbitales: planetas calientes, cálidos, templados, fríos o lejanos.
    \item Comparar si las redes física y orbital comparten estructura.
    \item Evaluar si la red conjunta preserva las relaciones observadas en los Mappers separados.
    \item Medir si la topología está sesgada por método de descubrimiento.
    \item Detectar planetas topológicamente raros: periféricos, puente, extremos o ambiguos.
    \item Explorar si las variables topológicas pueden usarse después como entradas para machine learning.
\end{enumerate}

\section{Hipótesis}

\subsection{Hipótesis principal}

La hipótesis central es:

\begin{quote}
La estructura físico-composicional y la estructura orbital de los exoplanetas no son equivalentes, pero están parcialmente acopladas en regiones específicas del espacio conjunto.
\end{quote}

En notación:

\begin{equation}
\mathcal{T}(X_{\phys}) \neq \mathcal{T}(X_{\orb}),
\end{equation}

pero:

\begin{equation}
\exists\, R\subseteq \mathcal{T}(X_{\joint}) \text{ tal que } R \text{ conecta poblaciones físicas y orbitales de forma estable.}
\end{equation}

\subsection{Hipótesis secundarias}

\begin{description}
    \item[H1: Transiciones físicas.] El Mapper físico mostrará ramas o regiones de baja densidad asociadas a transiciones entre planetas rocosos, super-Tierras y sub-Neptunos.
    \item[H2: Clases orbitales.] El Mapper orbital separará planetas por régimen energético y escala orbital, especialmente por período, insolación y temperatura de equilibrio.
    \item[H3: Acoplamiento parcial.] Algunas familias físicas estarán sobrerrepresentadas en ciertas regiones orbitales, pero no existirá una correspondencia uno-a-uno.
    \item[H4: Sesgo observacional.] Algunas ramas del Mapper estarán dominadas por método de descubrimiento, por ejemplo tránsito, velocidad radial, imagen directa o microlente.
    \item[H5: Valor predictivo.] Las características topológicas derivadas del Mapper mejorarán o explicarán modelos de machine learning frente a usar solamente variables tabulares originales.
\end{description}

\section{Datos y variables}

\subsection{Fuente de datos}

La fuente primaria será la tabla \texttt{PSCompPars} del NASA Exoplanet Archive. La tabla provee una vista más completa, con una fila por planeta confirmado, pero no necesariamente autoconsistente físicamente porque algunos valores pueden proceder de diferentes referencias o cálculos \cite{nasa_pscp_about,nasa_pscp_calc}. Por ello, se recomienda:

\begin{itemize}
    \item conservar referencias de masa, radio y densidad cuando sea posible;
    \item registrar proporción de valores faltantes por variable;
    \item distinguir valores empíricos de valores calculados;
    \item analizar sensibilidad a imputación;
    \item no tratar el catálogo como muestra uniforme del universo.
\end{itemize}

\subsection{Variables principales}

\begin{longtable}{p{0.17\linewidth}p{0.18\linewidth}p{0.17\linewidth}p{0.38\linewidth}}
\caption{Variables principales del proyecto.}\\
\toprule
\textbf{Bloque} & \textbf{Columna} & \textbf{Transformación} & \textbf{Justificación} \\
\midrule
\endfirsthead
\toprule
\textbf{Bloque} & \textbf{Columna} & \textbf{Transformación} & \textbf{Justificación} \\
\midrule
\endhead
Física & \texttt{pl\_rade} & $\log_{10}$ & Radio planetario en radios terrestres. Separa planetas pequeños, sub-Neptunos y gigantes. \\
Física & \texttt{pl\_bmasse} & $\log_{10}$ & Masa planetaria o mejor masa disponible en masas terrestres. Captura escala gravitacional. \\
Física & \texttt{pl\_dens} & $\log_{10}$ & Densidad planetaria. Aproxima composición: rocoso, volátil, gaseoso o extremo. \\
Orbital & \texttt{pl\_orbper} & $\log_{10}$ & Período orbital. Distingue planetas cercanos, cálidos y lejanos. \\
Orbital & \texttt{pl\_orbsmax} & $\log_{10}$ & Semieje mayor. Mide distancia orbital en unidades astronómicas. \\
Orbital-energética & \texttt{pl\_insol} & $\log_{10}$ & Insolación relativa. Relacionada con pérdida atmosférica, temperatura y habitabilidad aproximada. \\
Orbital-energética & \texttt{pl\_eqt} & robust scaler & Temperatura de equilibrio. Resume régimen térmico. \\
Dinámica opcional & \texttt{pl\_orbeccen} & sin log & Excentricidad. Captura circularidad y posibles historias dinámicas. \\
Sistema opcional & \texttt{sy\_pnum} & estandarizar & Número de planetas en el sistema. Útil para arquitectura. \\
Sistema opcional & \texttt{sy\_snum} & estandarizar & Número de estrellas del sistema. Útil para binariedad y contexto dinámico. \\
\bottomrule
\end{longtable}

\subsection{Variables de sesgo y control}

\begin{longtable}{p{0.21\linewidth}p{0.22\linewidth}p{0.46\linewidth}}
\caption{Variables observacionales para evaluar sesgo.}\\
\toprule
\textbf{Columna} & \textbf{Uso} & \textbf{Razón} \\
\midrule
\endfirsthead
\toprule
\textbf{Columna} & \textbf{Uso} & \textbf{Razón} \\
\midrule
\endhead
\texttt{discoverymethod} & etiqueta externa & Permite saber si un nodo, rama o comunidad está dominado por tránsito, velocidad radial, imagen directa, microlente u otro método. \\
\texttt{disc\_year} & gradiente o control & Captura sesgo temporal e historial instrumental. \\
\texttt{disc\_facility} & etiqueta externa & Puede revelar dominancia de una misión o instalación. \\
\texttt{sy\_dist} & control & La distancia al sistema afecta detectabilidad y calidad de caracterización. \\
\texttt{tran\_flag} & etiqueta binaria & Indica detección o caracterización por tránsito. \\
\texttt{rv\_flag} & etiqueta binaria & Indica señal de velocidad radial. \\
\texttt{ima\_flag} & etiqueta binaria & Indica imagen directa; suele asociarse con planetas masivos y lejanos. \\
\texttt{micro\_flag} & etiqueta binaria & Indica microlente; puede generar regiones con caracterización limitada. \\
\bottomrule
\end{longtable}

La regla metodológica será:

\begin{equation}
\text{entrada de Mapper} = \text{variables físicas u orbitales},
\end{equation}

\begin{equation}
\text{diagnóstico de sesgo} = \text{método de descubrimiento y controles observacionales}.
\end{equation}

\section{Construcción de los espacios de análisis}

\subsection{Espacio físico-composicional}

\begin{equation}
X_{\phys}=\left[\log_{10}\Rp,\log_{10}\Mp,\log_{10}\rhop\right].
\end{equation}

Este espacio busca capturar familias planetarias por tamaño, masa y composición.

\begin{equation}
(\Rp,\Mp,\rhop)\longrightarrow
\begin{cases}
\text{terrestres},\\
\text{super-Tierras},\\
\text{sub-Neptunos},\\
\text{gigantes gaseosos},\\
\text{super-Júpiteres}.
\end{cases}
\end{equation}

\subsection{Espacio orbital-energético}

\begin{equation}
X_{\orb}=\left[\log_{10}P,\log_{10}a,\log_{10}S,\Teq\right].
\end{equation}

Este espacio busca capturar clases orbitales y térmicas.

\begin{equation}
(P,a,S,\Teq)\longrightarrow
\begin{cases}
\text{planetas calientes},\\
\text{planetas cálidos},\\
\text{planetas templados},\\
\text{planetas fríos o lejanos}.
\end{cases}
\end{equation}

\subsection{Espacio conjunto}

\begin{equation}
X_{\joint}=\left[X_{\phys},X_{\orb}\right].
\end{equation}

Explícitamente:

\begin{equation}
X_{\joint}=\left[\log_{10}\Rp,\log_{10}\Mp,\log_{10}\rhop,\log_{10}P,\log_{10}a,\log_{10}S,\Teq\right].
\end{equation}

Este espacio permitirá evaluar si las estructuras separadas se integran en una misma topología.

\section{Preprocesamiento}

\subsection{Filtrado inicial}

Se recomienda iniciar con un subconjunto de planetas que tenga valores válidos para:

\begin{equation}
\{\texttt{pl\_rade},\texttt{pl\_bmasse},\texttt{pl\_dens},\texttt{pl\_orbper},\texttt{pl\_orbsmax},\texttt{pl\_insol},\texttt{pl\_eqt}\}.
\end{equation}

Luego se evaluará una versión ampliada con imputación.

\subsection{Transformación logarítmica}

Las variables positivas con varias órdenes de magnitud deben transformarse:

\begin{equation}
x' = \log_{10}(x).
\end{equation}

Esto aplica a:

\begin{equation}
\Rp,\Mp,\rhop,P,a,S,\texttt{sy\_dist}.
\end{equation}

\subsection{Escalamiento robusto}

Después del logaritmo se usará escalamiento robusto:

\begin{equation}
z_j=\frac{x_j-\operatorname{median}(x_j)}{\operatorname{IQR}(x_j)}.
\end{equation}

La razón es que las poblaciones exoplanetarias tienen colas largas y outliers reales.

\subsection{Manejo de valores faltantes}

Se recomiendan tres versiones:

\begin{enumerate}
    \item \textbf{Complete case:} usar solo planetas con todas las variables principales disponibles.
    \item \textbf{Imputación simple robusta:} imputar mediana por variable o por método de descubrimiento.
    \item \textbf{Imputación iterativa:} usar variables correlacionadas para predecir faltantes, reportando sensibilidad.
\end{enumerate}

La comparación entre versiones será parte del control de robustez.

\section{Metodología Mapper}

\subsection{Definición del algoritmo}

Dado un conjunto de datos $X$, una función filtro $f:X\to\R^k$, una cubierta solapada $\mathcal{U}$ del espacio del filtro y un algoritmo de clustering local $C$, Mapper produce un grafo:

\begin{equation}
G=\Mapper(X;f,\mathcal{U},C).
\end{equation}

Los pasos son:

\begin{enumerate}
    \item Aplicar un filtro o lens $f$ a los datos.
    \item Cubrir el espacio filtrado con intervalos o hipercubos solapados.
    \item Tomar preimágenes de cada elemento de la cubierta.
    \item Clusterizar localmente cada preimagen.
    \item Crear un nodo por cluster local.
    \item Conectar dos nodos si comparten al menos un punto original.
\end{enumerate}

\subsection{Filtros candidatos}

\begin{longtable}{p{0.22\linewidth}p{0.30\linewidth}p{0.38\linewidth}}
\caption{Filtros recomendados para Mapper.}\\
\toprule
\textbf{Filtro} & \textbf{Forma} & \textbf{Uso} \\
\midrule
\endfirsthead
\toprule
\textbf{Filtro} & \textbf{Forma} & \textbf{Uso} \\
\midrule
\endhead
PCA 1D & $f(x)=\operatorname{PC}_1(x)$ & Captura el gradiente dominante del espacio. \\
PCA 2D & $f(x)=(\operatorname{PC}_1,\operatorname{PC}_2)$ & Permite cubierta bidimensional y estructuras más ricas. \\
UMAP 2D & $f(x)=\operatorname{UMAP}(x)$ & Captura relaciones no lineales; útil como visualización, pero puede introducir artefactos. \\
Variable física & $f(x)=\log \Rp$ o $\log\Mp$ & Útil para estudiar transiciones por tamaño o masa. \\
Variable orbital & $f(x)=\log P$ o $\log S$ & Útil para ordenar regiones por escala orbital o energética. \\
Densidad local & $f(x)=\widehat{\rho}(x)$ & Resalta núcleos, periferias y outliers. \\
\bottomrule
\end{longtable}

\subsection{Clusterer local}

La opción inicial recomendada es DBSCAN o HDBSCAN, porque no obliga a fijar el número de clusters locales. Alternativamente, se puede usar Agglomerative Clustering o KMeans como baseline.

\section{Diseño experimental}

\subsection{Mappers principales}

\begin{equation}
G_{\phys}=\Mapper(X_{\phys};f_{\phys},\mathcal{U}_{\phys},C),
\end{equation}

\begin{equation}
G_{\orb}=\Mapper(X_{\orb};f_{\orb},\mathcal{U}_{\orb},C),
\end{equation}

\begin{equation}
G_{\joint}=\Mapper(X_{\joint};f_{\joint},\mathcal{U}_{\joint},C).
\end{equation}

\subsection{Parámetros a explorar}

\begin{longtable}{p{0.25\linewidth}p{0.25\linewidth}p{0.40\linewidth}}
\caption{Hiperparámetros de Mapper.}\\
\toprule
\textbf{Elemento} & \textbf{Rango inicial} & \textbf{Razón} \\
\midrule
\endfirsthead
\toprule
\textbf{Elemento} & \textbf{Rango inicial} & \textbf{Razón} \\
\midrule
\endhead
Número de intervalos & 8, 10, 12, 15, 20 & Controla resolución del grafo. \\
Overlap & 20\%, 30\%, 40\%, 50\% & Controla conectividad entre nodos. \\
Filtro & PCA 1D, PCA 2D, UMAP 2D, densidad & Evalúa estabilidad frente a lens. \\
Clusterer & DBSCAN, HDBSCAN, Agglomerative & Evalúa dependencia al algoritmo local. \\
Distancia & Euclidiana robusta, Mahalanobis opcional & Controla geometría del espacio. \\
\bottomrule
\end{longtable}

\section{Comparación de grafos}

\subsection{Descriptores globales}

Para cada grafo $G=(V,E)$ se calculará:

\begin{equation}
\phi(G)=\left[|V|,|E|,\beta_0,\beta_1,\operatorname{diam}(G),\bar d(G),C(G),Q(G)\right],
\end{equation}

donde:

\begin{equation}
\beta_1(G)=|E|-|V|+\beta_0(G).
\end{equation}

\begin{longtable}{p{0.22\linewidth}p{0.25\linewidth}p{0.43\linewidth}}
\caption{Métricas de red.}\\
\toprule
\textbf{Métrica} & \textbf{Símbolo} & \textbf{Interpretación} \\
\midrule
\endfirsthead
\toprule
\textbf{Métrica} & \textbf{Símbolo} & \textbf{Interpretación} \\
\midrule
\endhead
Nodos & $|V|$ & Resolución y fragmentación del Mapper. \\
Aristas & $|E|$ & Conectividad entre regiones locales. \\
Componentes & $\beta_0$ & Número de poblaciones separadas. \\
Ciclos & $\beta_1$ & Huecos o trayectorias cerradas en el grafo. \\
Diámetro & $\operatorname{diam}(G)$ & Extensión topológica máxima. \\
Camino promedio & $\bar d(G)$ & Separación media entre nodos. \\
Clustering coefficient & $C(G)$ & Tendencia a triángulos o cierre local. \\
Modularidad & $Q(G)$ & Fuerza de comunidades. \\
\bottomrule
\end{longtable}

\subsection{Distancia entre grafos por features}

La distancia simple será:

\begin{equation}
D_{\rm feat}(G_i,G_j)=\left\|\operatorname{zscore}(\phi(G_i))-\operatorname{zscore}(\phi(G_j))\right\|_2.
\end{equation}

\subsection{Distancia espectral}

Se calcularán autovalores del Laplaciano normalizado:

\begin{equation}
\mathcal{L}=I-D^{-1/2}AD^{-1/2}.
\end{equation}

La distancia espectral será:

\begin{equation}
D_{\rm spec}(G_i,G_j)=\left(\sum_{k=1}^{K}\left(\lambda_k^{(i)}-\lambda_k^{(j)}\right)^2\right)^{1/2}.
\end{equation}

\subsection{Distancia entre clases dentro del Mapper}

Si $A$ y $B$ son dos clases, por ejemplo super-Tierras y sub-Neptunos, definimos:

\begin{equation}
V_A=\{v\in V: v\text{ está enriquecido en clase }A\},
\end{equation}

\begin{equation}
V_B=\{v\in V: v\text{ está enriquecido en clase }B\}.
\end{equation}

La distancia topológica promedio entre clases será:

\begin{equation}
D_G(A,B)=\frac{1}{|V_A||V_B|}\sum_{u\in V_A}\sum_{v\in V_B}d_G(u,v),
\end{equation}

donde $d_G(u,v)$ es distancia de camino mínimo en el grafo.

\section{Análisis de sesgo por método de descubrimiento}

\subsection{Idea principal}

El método de descubrimiento no se usará para construir el Mapper principal. Se usará como etiqueta externa para estudiar si las regiones topológicas están dominadas por una técnica observacional.

\begin{equation}
G=\Mapper(X_{\phys}\text{ o }X_{\orb}\text{ o }X_{\joint}),
\end{equation}

\begin{equation}
\texttt{discoverymethod}\notin X,\qquad \texttt{discoverymethod}\in \text{etiquetas externas}.
\end{equation}

\subsection{Distribución local por nodo}

Para cada nodo $v$ del Mapper, sea $C_v$ el conjunto de planetas contenidos en él. Para cada método $m$:

\begin{equation}
p_v(m)=\frac{|\{x_i\in C_v:\operatorname{method}(x_i)=m\}|}{|C_v|}.
\end{equation}

El método dominante es:

\begin{equation}
m^*(v)=\arg\max_m p_v(m).
\end{equation}

La pureza local es:

\begin{equation}
\operatorname{purity}(v)=\max_m p_v(m).
\end{equation}

\subsection{Entropía e índice de sesgo}

La entropía por método en un nodo es:

\begin{equation}
H(v)=-\sum_{m=1}^{M}p_v(m)\log p_v(m).
\end{equation}

El índice normalizado de sesgo será:

\begin{equation}
B(v)=1-\frac{H(v)}{\log M}.
\end{equation}

\begin{longtable}{p{0.22\linewidth}p{0.62\linewidth}}
\caption{Interpretación del índice de sesgo $B(v)$.}\\
\toprule
\textbf{Valor} & \textbf{Interpretación} \\
\midrule
\endfirsthead
\toprule
\textbf{Valor} & \textbf{Interpretación} \\
\midrule
\endhead
$B(v)\approx 0$ & Nodo mezclado; baja dominancia de método. \\
$B(v)\approx 1$ & Nodo dominado por un solo método. \\
$B(v)>0.7$ & Sesgo local fuerte. \\
$B(v)>0.9$ & Sesgo local extremo. \\
\bottomrule
\end{longtable}

\subsection{Enriquecimiento por método}

Para saber hacia cuál método está sesgada una región, compararemos la proporción local contra la proporción global:

\begin{equation}
E(v,m)=\log_2\left(\frac{p_v(m)+\epsilon}{p_{\rm global}(m)+\epsilon}\right).
\end{equation}

\begin{longtable}{p{0.22\linewidth}p{0.62\linewidth}}
\caption{Interpretación del enriquecimiento $E(v,m)$.}\\
\toprule
\textbf{Valor} & \textbf{Interpretación} \\
\midrule
\endfirsthead
\toprule
\textbf{Valor} & \textbf{Interpretación} \\
\midrule
\endhead
$E(v,m)>0$ & Método sobrerrepresentado localmente. \\
$E(v,m)=0$ & Método en proporción esperada. \\
$E(v,m)<0$ & Método subrepresentado. \\
$E(v,m)>1$ & Método al menos dos veces más frecuente que globalmente. \\
$E(v,m)>2$ & Método al menos cuatro veces más frecuente que globalmente. \\
\bottomrule
\end{longtable}

\subsection{Asociación global}

Se construirá una tabla de contingencia:

\begin{equation}
N_{v,m}=\text{número de planetas en nodo }v\text{ detectados por método }m.
\end{equation}

Se aplicará una prueba $\chi^2$ de asociación entre nodo Mapper y método de descubrimiento. Como tamaño de efecto se usará $V$ de Cramer:

\begin{equation}
V_{\rm Cramer}=\sqrt{\frac{\chi^2}{n\min(r-1,c-1)}}.
\end{equation}

\subsection{Prueba por permutación}

Para validar si el sesgo es mayor que azar, se permutarán las etiquetas de método:

\begin{equation}
\operatorname{method}(x_i)\mapsto \operatorname{method}(x_{\pi(i)}).
\end{equation}

El valor $p$ empírico será:

\begin{equation}
p=\frac{1+\sum_{k=1}^{K}\mathbf{1}\left[B^{(k)}\geq B_{\rm obs}\right]}{K+1}.
\end{equation}

\section{Validación y robustez}

\subsection{Bootstrap}

Se generarán muestras bootstrap:

\begin{equation}
X^{(b)}\sim \operatorname{Bootstrap}(X),\qquad b=1,\dots,B.
\end{equation}

Para cada muestra se calcularán:

\begin{equation}
G_{\phys}^{(b)},\quad G_{\orb}^{(b)},\quad G_{\joint}^{(b)}.
\end{equation}

Una estructura será considerada estable si aparece en una fracción alta de réplicas.

\subsection{Grid de hiperparámetros}

Se repetirá Mapper sobre una malla:

\begin{equation}
\Theta=\{(n_{\rm intervals},\operatorname{overlap},f,C)\}.
\end{equation}

La estabilidad se evaluará con:

\begin{equation}
\Delta_{ij}=D(G_{\theta_i},G_{\theta_j}).
\end{equation}

Si una conclusión depende de un solo valor de $\theta$, se considera frágil.

\subsection{Control nulo físico-orbital}

Para probar acoplamiento entre física y órbita, se rompe la correspondencia planeta-a-planeta:

\begin{equation}
(X_{\phys}^{(i)},X_{\orb}^{(i)})\mapsto (X_{\phys}^{(i)},X_{\orb}^{(\pi(i))}).
\end{equation}

Si el Mapper conjunto real muestra estructura más fuerte que el permutado, hay evidencia de asociación físico-orbital no trivial.

\section{Extensión hacia machine learning}

\subsection{Preguntas supervisadas}

El proyecto puede escalar hacia machine learning con tres tareas:

\begin{align}
f_1 &: X_{\orb}\to y_{\phys},\\
f_2 &: X_{\phys}\to y_{\orb},\\
f_3 &: X_{\joint}\to y_{\rm method}.
\end{align}

Donde:

\begin{itemize}
    \item $y_{\phys}$: familia físico-composicional.
    \item $y_{\orb}$: clase orbital-energética.
    \item $y_{\rm method}$: método de descubrimiento.
\end{itemize}

\subsection{Features topológicas para ML}

A cada planeta se le pueden agregar características del nodo Mapper donde aparece:

\begin{longtable}{p{0.26\linewidth}p{0.58\linewidth}}
\caption{Variables topológicas derivadas.}\\
\toprule
\textbf{Feature} & \textbf{Significado} \\
\midrule
\endfirsthead
\toprule
\textbf{Feature} & \textbf{Significado} \\
\midrule
\endhead
\texttt{mapper\_node\_id} & Identificador del nodo Mapper. \\
\texttt{node\_size} & Número de planetas en el nodo; aproxima densidad local. \\
\texttt{node\_degree} & Conectividad del nodo. \\
\texttt{node\_betweenness} & Capacidad del nodo de actuar como puente entre regiones. \\
\texttt{node\_closeness} & Centralidad global. \\
\texttt{node\_community} & Comunidad topológica. \\
\texttt{node\_bias\_score} & Sesgo por método en el nodo. \\
\texttt{node\_class\_entropy} & Mezcla de clases físicas u orbitales. \\
\texttt{distance\_to\_core} & Distancia a una región central o densa. \\
\bottomrule
\end{longtable}

Entonces se compara:

\begin{equation}
\operatorname{score}(X)\quad \text{vs.}\quad \operatorname{score}([X,\phi_{\rm Mapper}(x)]).
\end{equation}

Si el segundo mejora, la topología aporta información predictiva.

\subsection{Modelos recomendados}

\begin{longtable}{p{0.24\linewidth}p{0.32\linewidth}p{0.34\linewidth}}
\caption{Modelos de machine learning recomendados.}\\
\toprule
\textbf{Modelo} & \textbf{Uso} & \textbf{Razón} \\
\midrule
\endfirsthead
\toprule
\textbf{Modelo} & \textbf{Uso} & \textbf{Razón} \\
\midrule
\endhead
Logistic Regression & baseline supervisado & Interpretable, útil para comprobar señal mínima. \\
Random Forest & clasificación y sesgo & Captura no linealidades e importancia de variables. \\
XGBoost/LightGBM & rendimiento tabular & Fuerte para datos heterogéneos y faltantes tratados. \\
SVM & separación no lineal & Útil con espacios reducidos y clases complejas. \\
Isolation Forest & anomalías & Detecta planetas raros en $X_{\joint}$. \\
Node2Vec & embeddings de nodos & Aprende representación de nodos del Mapper. \\
Graph Neural Networks & extensión avanzada & Usar solo si hay suficiente justificación y múltiples grafos. \\
\bottomrule
\end{longtable}

\section{Resultados esperados}

\subsection{Resultado científico ideal}

Una conclusión fuerte tendría la forma:

\begin{quote}
La topología física y la topología orbital no son idénticas, pero muestran acoplamientos locales. El Mapper conjunto revela regiones puente entre poblaciones físico-orbitales, algunas consistentes con estructuras conocidas de poblaciones exoplanetarias y otras dominadas por sesgo de método de descubrimiento.
\end{quote}

\subsection{Resultados alternativos valiosos}

\begin{longtable}{p{0.30\linewidth}p{0.58\linewidth}}
\caption{Escenarios posibles y su interpretación.}\\
\toprule
\textbf{Resultado} & \textbf{Interpretación} \\
\midrule
\endfirsthead
\toprule
\textbf{Resultado} & \textbf{Interpretación} \\
\midrule
\endhead
$G_{\phys}\approx G_{\orb}$ & Fuerte acoplamiento entre composición y órbita; debe validarse contra sesgo. \\
$G_{\phys}\not\approx G_{\orb}$ & La física y la órbita organizan los datos de forma distinta; resultado no fallido. \\
$G_{\joint}$ preserva puentes & Evidencia de integración físico-orbital. \\
$G_{\joint}$ dominado por método & La topología refleja función de selección o detectabilidad. \\
Alta estabilidad bootstrap & Conclusión robusta. \\
Baja estabilidad bootstrap & La estructura depende de muestra o parámetros. \\
Anomalías estables & Lista de planetas científicamente interesantes para análisis posterior. \\
\bottomrule
\end{longtable}

\section{Plan de trabajo}

\begin{longtable}{p{0.12\linewidth}p{0.30\linewidth}p{0.45\linewidth}}
\caption{Cronograma conceptual.}\\
\toprule
\textbf{Fase} & \textbf{Objetivo} & \textbf{Entregable} \\
\midrule
\endfirsthead
\toprule
\textbf{Fase} & \textbf{Objetivo} & \textbf{Entregable} \\
\midrule
\endhead
1 & Descarga y documentación de datos & Dataset crudo, diccionario de variables, resumen de faltantes. \\
2 & Limpieza y preprocesamiento & Dataset procesado con logs, escalamiento e imputación controlada. \\
3 & Mapper físico & Grafo $G_{\phys}$, métricas y visualizaciones. \\
4 & Mapper orbital & Grafo $G_{\orb}$, métricas y visualizaciones. \\
5 & Mapper conjunto & Grafo $G_{\joint}$ y comparación con los dos anteriores. \\
6 & Sesgo por método & Tablas de enriquecimiento, entropía, Cramer $V$ y permutación. \\
7 & Robustez & Bootstrap, grid de hiperparámetros y estabilidad. \\
8 & Anomalías & Catálogo de planetas puente, periféricos y raros. \\
9 & ML exploratorio & Modelos baseline y comparación con features topológicas. \\
10 & Reporte final & Atlas, conclusiones, limitaciones y apéndice reproducible. \\
\bottomrule
\end{longtable}

\section{Herramientas computacionales}

\begin{longtable}{p{0.22\linewidth}p{0.28\linewidth}p{0.38\linewidth}}
\caption{Herramientas propuestas.}\\
\toprule
\textbf{Categoría} & \textbf{Herramientas} & \textbf{Uso} \\
\midrule
\endfirsthead
\toprule
\textbf{Categoría} & \textbf{Herramientas} & \textbf{Uso} \\
\midrule
\endhead
Datos & NASA Exoplanet Archive TAP, astroquery & Descarga y consulta de \texttt{PSCompPars}. \\
Procesamiento & pandas, numpy & Limpieza, transformación, imputación, tablas. \\
Escalamiento & scikit-learn & RobustScaler, pipelines, imputación, modelos. \\
Mapper & KeplerMapper, giotto-tda, custom NetworkX pipeline & Construcción de grafos Mapper. \\
Grafos & NetworkX, python-igraph & Métricas, comunidades, centralidades. \\
Topología & giotto-tda, ripser, persim & Persistencia y diagramas topológicos si se incorporan. \\
Visualización & matplotlib, plotly, pyvis & Gráficas, redes interactivas, reportes. \\
ML & scikit-learn, XGBoost/LightGBM & Modelos supervisados e interpretabilidad inicial. \\
Interpretabilidad & permutation importance, SHAP & Variables que explican método o clase. \\
Reproducibilidad & notebooks, scripts, conda/poetry, Git & Pipeline reproducible. \\
\bottomrule
\end{longtable}

\section{Pseudocódigo del pipeline}

\begin{lstlisting}[language=Python,caption={Pseudocódigo general del proyecto.}]
# 1. Cargar datos
raw = load_pscomppars()

# 2. Seleccionar variables
features_phys = ["pl_rade", "pl_bmasse", "pl_dens"]
features_orb  = ["pl_orbper", "pl_orbsmax", "pl_insol", "pl_eqt"]
features_joint = features_phys + features_orb
bias_vars = ["discoverymethod", "disc_year", "sy_dist",
             "tran_flag", "rv_flag", "ima_flag", "micro_flag"]

# 3. Preprocesar
X_phys  = preprocess(raw[features_phys], log_cols=features_phys)
X_orb   = preprocess(raw[features_orb], log_cols=["pl_orbper", "pl_orbsmax", "pl_insol"])
X_joint = concatenate([X_phys, X_orb])

# 4. Construir Mapper
G_phys  = build_mapper(X_phys, lens="pca2", cover_grid=theta)
G_orb   = build_mapper(X_orb, lens="pca2", cover_grid=theta)
G_joint = build_mapper(X_joint, lens="pca2", cover_grid=theta)

# 5. Medir grafos
metrics_phys  = graph_metrics(G_phys)
metrics_orb   = graph_metrics(G_orb)
metrics_joint = graph_metrics(G_joint)

# 6. Comparar grafos
D_phys_orb   = graph_distance(G_phys, G_orb)
D_phys_joint = graph_distance(G_phys, G_joint)
D_orb_joint  = graph_distance(G_orb, G_joint)

# 7. Diagnosticar sesgo por metodo
bias_table = node_method_enrichment(G_joint, labels=raw["discoverymethod"])
cramer_v = cramer_association(G_joint.node_memberships, raw["discoverymethod"])
p_value = permutation_test_method_bias(G_joint, raw["discoverymethod"])

# 8. Robustez
stability = bootstrap_mapper_stability(X_joint, labels=raw["discoverymethod"])

# 9. Anomalias
anomalies = topological_anomaly_scores(G_joint)

# 10. ML exploratorio
model_method = train_classifier(X_joint, y=raw["discoverymethod"])
model_topo = train_classifier(add_mapper_features(X_joint, G_joint), y=target)
\end{lstlisting}

\section{Riesgos y limitaciones}

\begin{longtable}{p{0.24\linewidth}p{0.36\linewidth}p{0.28\linewidth}}
\caption{Riesgos metodológicos.}\\
\toprule
\textbf{Riesgo} & \textbf{Consecuencia} & \textbf{Mitigación} \\
\midrule
\endfirsthead
\toprule
\textbf{Riesgo} & \textbf{Consecuencia} & \textbf{Mitigación} \\
\midrule
\endhead
Sesgo de detección & Ramas reflejan telescopios, no planetas & Etiquetar por método, enriquecimiento, permutación. \\
Faltantes no aleatorios & La imputación puede crear estructura falsa & Comparar complete case vs imputado. \\
Variables redundantes & Algunas variables derivan de otras & Reportar correlaciones y análisis de sensibilidad. \\
Sensibilidad de Mapper & Cambia por lens, cover o clusterer & Grid de hiperparámetros y bootstrap. \\
Interpretación visual excesiva & Conclusiones subjetivas & Métricas de grafos y pruebas estadísticas. \\
Muestra no uniforme & No inferir ocurrencia universal directamente & Hablar de catálogo observado, no de universo subyacente completo. \\
\bottomrule
\end{longtable}

\section{Conclusión propuesta}

El proyecto es viable y puede ser impactante si se formula como un atlas topológico validado. La contribución no debe ser únicamente producir visualizaciones, sino construir un marco para separar tres componentes:

\begin{equation}
\text{estructura observada}=\text{estructura astrofísica}+\text{sesgo observacional}+\text{artefacto metodológico}.
\end{equation}

La salida final debe permitir afirmar, con evidencia cuantitativa, cuáles regiones del espacio exoplanetario son:

\begin{itemize}
    \item físicamente interpretables;
    \item orbitalmente estructuradas;
    \item parcialmente acopladas entre física y órbita;
    \item dominadas por método de descubrimiento;
    \item topológicamente raras o candidatas a análisis posterior.
\end{itemize}

La tesis operativa del proyecto puede escribirse así:

\begin{quote}
Construiremos un atlas topológico de exoplanetas usando Mapper para comparar la estructura física, orbital y conjunta de la población observada, incorporando métricas de grafos, estabilidad por bootstrap, controles por método de descubrimiento y una ruta hacia machine learning interpretable.
\end{quote}

\section*{Recapitulación simbólica}
\addcontentsline{toc}{section}{Recapitulación simbólica}

\begin{equation}
\boxed{
X_{\phys}=[\log\Rp,\log\Mp,\log\rhop],\quad
X_{\orb}=[\log P,\log a,\log S,\Teq],\quad
X_{\joint}=[X_{\phys},X_{\orb}]
}
\end{equation}

\begin{equation}
\boxed{
G_{\phys}=\Mapper(X_{\phys}),\quad
G_{\orb}=\Mapper(X_{\orb}),\quad
G_{\joint}=\Mapper(X_{\joint})
}
\end{equation}

\begin{equation}
\boxed{
\text{Conclusión fuerte}\iff
\text{topología interpretable}+
\text{estabilidad}+
\text{control de sesgo}+
\text{validación estadística}
}
\end{equation}

\newpage
\begin{thebibliography}{9}

\bibitem{singh2007mapper}
Gurjeet Singh, Facundo M\'emoli, and Gunnar Carlsson.
\newblock \emph{Topological Methods for the Analysis of High Dimensional Data Sets and 3D Object Recognition}.
\newblock Eurographics Symposium on Point-Based Graphics, 2007.
\newblock \url{https://research.math.osu.edu/tgda/mapperPBG.pdf}

\bibitem{nasa_pscp_about}
NASA Exoplanet Archive.
\newblock \emph{About the Planetary Systems Composite Planet Data Table}.
\newblock 2023. Accessed 2026-04-26.
\newblock \url{https://exoplanetarchive.ipac.caltech.edu/docs/pscp_about.html}

\bibitem{nasa_ps_columns}
NASA Exoplanet Archive.
\newblock \emph{Planetary Systems and Planetary Systems Composite Parameters Table Definitions}.
\newblock 2025. Accessed 2026-04-26.
\newblock \url{https://exoplanetarchive.ipac.caltech.edu/docs/API_PS_columns.html}

\bibitem{nasa_pscp_calc}
NASA Exoplanet Archive.
\newblock \emph{How the Archive Calculates Values in the Planetary Systems Composite Parameters Table}.
\newblock 2024. Accessed 2026-04-26.
\newblock \url{https://exoplanetarchive.ipac.caltech.edu/docs/pscp_calc.html}

\bibitem{fulton2017gap}
Benjamin J. Fulton et al.
\newblock \emph{The California-Kepler Survey. III. A Gap in the Radius Distribution of Small Planets}.
\newblock The Astronomical Journal, 154(3):109, 2017.
\newblock \url{https://ui.adsabs.harvard.edu/abs/2017AJ....154..109F/abstract}

\bibitem{mazeh2016desert}
Tsevi Mazeh, Tomer Holczer, and Simchon Faigler.
\newblock \emph{Dearth of short-period Neptunian exoplanets: a desert in period-mass and period-radius planes}.
\newblock Astronomy \& Astrophysics, 589:A75, 2016.
\newblock \url{https://www.aanda.org/articles/aa/abs/2016/05/aa28065-15/aa28065-15.html}

\end{thebibliography}

\end{document}
